% Źródła: Hartshorne s. 304--311, 318 (§34, zad. 34.13--34.14); Greenberg 2010, s. 208.

\section{Trzy typy płaszczyzn} % lang-pl
\section{Tre tipi di piani} % lang-it
\label{sekcja_trzy_typy}

Saccheri dowodził swoich trzech hipotez ciągłością, a my wiemy dziś, że ciągłość jest zbędna. % lang-pl
W tej sekcji zobaczymy, jak jego wyniki wyglądają w dowolnej płaszczyźnie Hilberta, jak nazwiemy trzy przypadki i jak defekt trójkąta zacznie mierzyć odstępstwo od geometrii euklidesowej. % lang-pl
Saccheri dimostrava le sue tre ipotesi con la continuità, e oggi sappiamo che la continuità è superflua. % lang-it
In questa sezione vedremo come appaiono i suoi risultati in un piano di Hilbert qualsiasi, come chiameremo i tre casi e come il difetto di un triangolo comincerà a misurare lo scarto dalla geometria euclidea. % lang-it

Hartshorne \cite[s. 307]{hartshorne_2000} zauważa, że Saccheri odróżni hipotezę kąta ostrego, prostego i rozwartego według tego, czy kąty przy górnej podstawie są ostre, proste czy rozwarte; udowodni, że jeśli któraś zachodzi dla jednego czworokąta, to dla wszystkich, ale użyje ciągłości (w postaci twierdzenia o wartości pośredniej). % lang-pl
Hartshorne \cite[p. 307]{hartshorne_2000} osserva che Saccheri distinguerà l'ipotesi dell'angolo acuto, retto e ottuso a seconda che gli angoli alla base superiore siano acuti, retti o ottusi; dimostrerà che se una vale per un quadrilatero, vale per tutti, ma userà la continuità (nella forma del teorema del valore intermedio). % lang-it

\begin{theorem}
[Saccheriego] % lang-pl
[di Saccheri] % lang-it
\index{twierdzenie!Saccheriego (o trzech hipotezach)} % lang-pl
\index{teorema!di Saccheri (sulle tre ipotesi)} % lang-it
	W płaszczyźnie Hilberta, jeśli jeden czworokąt Saccheriego ma kąty przy górnej podstawie ostre, to wszystkie mają ostre; to samo dla kątów prostych i rozwartych. % lang-pl
	In un piano di Hilbert, se un quadrilatero di Saccheri ha gli angoli alla base superiore acuti, li hanno acuti tutti; lo stesso vale per angoli retti e ottusi. % lang-it
\end{theorem}

Dowód (bez ciągłości, przez porównywanie czworokątów o równych środkowych): Hartshorne \cite[s. 307--309]{hartshorne_2000} (tw. 34.5, poprzedzone prop. 34.2--34.4). % lang-pl
Dimostrazione (senza continuità, confrontando quadrilateri con la stessa mediana): Hartshorne \cite[p. 307--309]{hartshorne_2000} (teor. 34.5, preceduto dalle prop. 34.2--34.4). % lang-it

\begin{theorem}
	W dowolnej płaszczyźnie Hilberta zachodzi jedna z trzech możliwości. % lang-pl
	Jeśli istnieje trójkąt o sumie kątów mniejszej od dwóch kątów prostych, to każdy trójkąt ma sumę mniejszą. % lang-pl
	Następujące warunki są równoważne: istnieje trójkąt o sumie kątów równej dwóm kątom prostym; istnieje prostokąt; każdy trójkąt ma sumę kątów równą dwóm kątom prostym. % lang-pl
	Jeśli istnieje trójkąt o sumie większej od dwóch kątów prostych, to każdy trójkąt ma sumę większą. % lang-pl
	In un piano di Hilbert vale una delle tre possibilità. % lang-it
	Se esiste un triangolo con somma degli angoli minore di due angoli retti, allora ogni triangolo ha somma minore. % lang-it
	Le seguenti condizioni sono equivalenti: esiste un triangolo con somma degli angoli uguale a due angoli retti; esiste un rettangolo; ogni triangolo ha somma degli angoli uguale a due angoli retti. % lang-it
	Se esiste un triangolo con somma maggiore di due angoli retti, allora ogni triangolo ha somma maggiore. % lang-it
\end{theorem}

Dowód: Hartshorne \cite[s. 309--311]{hartshorne_2000} (tw. 34.7), na podstawie prop. 34.6: dla każdego trójkąta istnieje czworokąt Saccheriego, którego suma dwóch kątów przy górnej podstawie jest równa sumie kątów trójkąta \cite[s. 310]{hartshorne_2000}. % lang-pl
Dimostrazione: Hartshorne \cite[p. 309--311]{hartshorne_2000} (teor. 34.7), sulla base della prop. 34.6: per ogni triangolo esiste un quadrilatero di Saccheri la cui somma dei due angoli alla base superiore è uguale alla somma degli angoli del triangolo \cite[p. 310]{hartshorne_2000}. % lang-it

Te trzy przypadki Hartshorne nazywa geometrią \emph{półhiperboliczną} (suma mniejsza od $\pi$), \emph{półeuklidesową} (równa $\pi$) i \emph{półeliptyczną} (większa od $\pi$); słowo „hiperboliczna'' rezerwuje dla geometrii spełniającej aksjomat Hilberta o równoległych granicznych (zobacz sekcję \ref{sekcja_hiperboliczna}), bo półhiperboliczne płaszczyzny bez tego aksjomatu też istnieją \cite[s. 311]{hartshorne_2000}. % lang-pl
Przypadek półeliptyczny, w którym Saccheri odrzucił hipotezę kąta rozwartego, odkryje w 1900 roku Dehn, nazywając takie geometrie nielegendrowskimi; przykład to mały kawałek sfery nad ciałem nieArchimedesowym \cite[s. 311, 318]{hartshorne_2000}. % lang-pl
Hartshorne chiama questi tre casi geometria \emph{semiiperbolica} (somma minore di $\pi$), \emph{semieuclidea} (uguale a $\pi$) e \emph{semiellittica} (maggiore di $\pi$); la parola «iperbolica» la riserva alle geometrie che soddisfano l'assioma di Hilbert sulle parallele limite (si veda la sezione \ref{sekcja_hiperboliczna}), perché esistono anche piani semiiperbolici privi di tale assioma \cite[p. 311]{hartshorne_2000}. % lang-it
Il caso semiellittico, in cui Saccheri scartò l'ipotesi dell'angolo ottuso, lo scoprirà nel 1900 Dehn, chiamando tali geometrie non legendriane; un esempio è un piccolo pezzo di sfera su un corpo non archimedeo \cite[p. 311, 318]{hartshorne_2000}. % lang-it
\index[persons]{Dehn, Max}%
\index{geometria!półhiperboliczna} % lang-pl
\index{geometria!semiiperbolica} % lang-it

\begin{definition}
[defekt trójkąta] % lang-pl
[difetto di un triangolo] % lang-it
	Defektem trójkąta nazywamy liczbę $\delta = 2\mathrm{RA} - (\alpha + \beta + \gamma)$, gdzie RA oznacza kąt prosty; jest ona zerem dla trójkąta euklidesowego, dodatnia w płaszczyźnie półhiperbolicznej i ujemna w półeliptycznej. % lang-pl
	Il difetto di un triangolo è il numero $\delta = 2\mathrm{RA} - (\alpha + \beta + \gamma)$, dove RA indica l'angolo retto; è zero per un triangolo euclideo, positivo in un piano semiiperbolico e negativo in uno semiellittico. % lang-it
\end{definition}

Hartshorne \cite[s. 311]{hartshorne_2000}. % lang-pl
Hartshorne \cite[p. 311]{hartshorne_2000}. % lang-it

\begin{proposition}
	Jeśli odcinek z wierzchołka dzieli trójkąt na dwa trójkąty, to defekt całości jest sumą defektów części: $\delta(ABC) = \delta(ABD) + \delta(BCD)$. % lang-pl
	Se un segmento uscente da un vertice divide un triangolo in due triangoli, il difetto del tutto è la somma dei difetti delle parti: $\delta(ABC) = \delta(ABD) + \delta(BCD)$. % lang-it
\end{proposition}

Dowód: Hartshorne \cite[s. 311--312]{hartshorne_2000} (lemat 34.8). % lang-pl
Dimostrazione: Hartshorne \cite[p. 311--312]{hartshorne_2000} (lemma 34.8). % lang-it

Z tego wprost wyniknie pole trójkąta jako defekt (zobacz sekcję \ref{sekcja_suma_katow}), a w sekcji o polu nieeuklidesowym Hartshorne \cite[s. 326 nn.]{hartshorne_2000} rozwija to systematycznie. % lang-pl
Da qui discenderà direttamente l'area del triangolo come difetto (si veda la sezione \ref{sekcja_suma_katow}), e nella sezione sull'area non euclidea Hartshorne \cite[p. 326 e segg.]{hartshorne_2000} lo sviluppa sistematicamente. % lang-it

Zadania: czworokąt Lamberta w trzech geometriach (zad. 34.2), kąt wpisany w półokrąg (zad. 34.3) i przystawanie trójkątów o równych kątach (zad. 34.4) -- Hartshorne \cite[s. 316]{hartshorne_2000}. % lang-pl
Esercizi: il quadrilatero di Lambert nelle tre geometrie (es. 34.2), l'angolo inscritto in una semicirconferenza (es. 34.3) e la congruenza di triangoli con angoli uguali (es. 34.4) -- Hartshorne \cite[p. 316]{hartshorne_2000}. % lang-it
